\documentclass[11pt]{article}

% ---------- packages ----------
\usepackage[top=0.75in,bottom=0.75in,left=0.9in,right=0.9in,headheight=14pt]{geometry}
\usepackage{amsmath,amssymb}
\usepackage{enumitem}
\usepackage{fancyhdr}
\usepackage{lastpage}
\usepackage{needspace}
\usepackage{hyperref}
\hypersetup{colorlinks=true,urlcolor=black,linkcolor=black}

% ---------- header / footer ----------
\pagestyle{fancy}
\fancyhf{}
\lhead{CS 418: Intro to Data Science}
\rhead{UIC, Fall 2026}
\cfoot{\thepage}
\renewcommand{\headrulewidth}{0.4pt}

% ---------- worksheet metadata ----------
\newcommand{\worksheetnumber}{1}
\newcommand{\topictitle}{Intro, Stats Review, and the Data Science Lifecycle}

% ---------- problem command ----------
\newcounter{problem}
\newcommand{\problem}[1]{%
  \refstepcounter{problem}%
  \vspace{0.55em}%
  \noindent\textbf{Problem \theproblem.} #1\par
}
\newcommand{\hint}[1]{\vspace{0.3em}\noindent\textit{Hint: #1}\par}
% Global scale for answer space: change 3/5 to tighten or loosen every gap at once.
\newcommand{\worklines}[1]{\vspace{\dimexpr#1*1/2\relax}}

\setlist[enumerate,1]{label=(\alph*), itemsep=0.5em, topsep=0.2em, parsep=0pt}

\begin{document}

% ================= HEADER BLOCK =================
\begin{center}
  {\large\textbf{Worksheet \#\worksheetnumber}}\\[0.25em]
  {\large \topictitle}
\end{center}

\vspace{0.5em}
\noindent
\begin{tabular}{@{}l@{}}
Name: \rule{2.2in}{0.4pt}\\[0.55em]
NetID: \rule{2.2in}{0.4pt}\\[0.55em]
Date: \rule{2.2in}{0.4pt}
\end{tabular}

\vspace{0.9em}

\noindent\textbf{Learning objectives:}
\begin{itemize}[label=\textbullet, leftmargin=1.6em, itemsep=0.1em, topsep=0.2em]
  \item Describe data science as an interdisciplinary field and explain why practitioners define it differently.
  \item Compute classical, complementary, and conditional probabilities, and apply Bayes' theorem.
  \item Trace a real question through the five stages of the data science lifecycle.
\end{itemize}

\vspace{0.8em}
\hrule
\vspace{0.5em}

% ================= PROBLEMS =================

\problem{Drew Conway's Venn diagram places data science at the intersection of hacking skills, math and statistics knowledge, and substantive expertise.}
\begin{enumerate}
  \item Which two circles overlap to form the ``danger zone''? Explain in one or two sentences why that combination is dangerous.
  \worklines{2.5em}
  \item We described data science as an \emph{Elephantidae} --- a family of related methods. Give an example of two people whose jobs both count as data science but who would describe the field differently.
  \worklines{3em}
\end{enumerate}

\problem{A bag contains 6 blue, 3 red, and 5 yellow marbles. One marble is drawn at random.}
\begin{enumerate}
  \item How many outcomes are in the sample space $S$?
  \worklines{2em}
  \item Let $E$ be the event that the marble is blue or red. Compute $P(E) = \dfrac{|E|}{|S|}$.
  \worklines{2.5em}
  \item Let $Y$ be the event that the marble is yellow. Compute $P(Y)$ and verify that $E$ and $Y$ are complementary events.
  \worklines{2.5em}
  \item Now two marbles are drawn \emph{without replacement}. What is the probability that both are blue?
  \worklines{3em}
\end{enumerate}

\problem{Match each term to the correct description by writing the letter in the blank.}

\vspace{0.5em}
\noindent
\begin{tabular}{@{}p{0.32\textwidth}p{0.62\textwidth}@{}}
\rule{0.4in}{0.4pt} Experiment & (A) The set of all possible results \\[0.6em]
\rule{0.4in}{0.4pt} Outcome & (B) A set of one or more results we care about \\[0.6em]
\rule{0.4in}{0.4pt} Event & (C) A process or action with an uncertain result \\[0.6em]
\rule{0.4in}{0.4pt} Sample space & (D) Two events where one occurs if and only if the other does not \\[0.6em]
\rule{0.4in}{0.4pt} Complementary events & (E) A single possible result of the process \\
\end{tabular}

\needspace{7\baselineskip}
\problem{Two of 13 distinct cards labeled 1--13 are drawn \emph{without replacement}. Let $A$ = both cards odd, and $B$ = the sum is even.}
\begin{enumerate}
  \item Explain in one sentence why the sum is even exactly when both cards are odd or both cards are even.
  \worklines{2em}
  \item Count $n(B)$ using binomial coefficients.
  \worklines{3em}
  \item Count $n(A \cap B)$.
  \worklines{2.5em}
  \item Compute $P(A \mid B) = \dfrac{n(A \cap B)}{n(B)}$ and simplify.
  \worklines{3em}
\end{enumerate}
\hint{There are 7 odd labels and 6 even labels. Order does not matter, so count unordered pairs.}

\problem{A fraud-detection model flags transactions for review. $0.5\%$ of transactions are fraudulent; the model flags $95\%$ of fraudulent transactions and $2\%$ of legitimate ones. Let $F$ = fraudulent and $R$ = flagged.}
\begin{enumerate}
  \item Write down $P(F)$, $P(R \mid F)$, and $P(R \mid F')$ from the description.
  \worklines{2.5em}
  \item Use the law of total probability to compute $P(R)$:
  \[ P(R) = P(R \mid F)P(F) + P(R \mid F')P(F') \]
  \worklines{3em}
  \item Apply Bayes' theorem to compute $P(F \mid R)$, the probability that a flagged transaction is actually fraudulent.
  \worklines{3.5em}
\end{enumerate}
\hint{Think about how many legitimate transactions there are compared to fraudulent ones.}

\problem{\textbf{Scenario:} UIC wants to know whether the new shuttle schedule reduced student wait times on campus. For each stage of the lifecycle, write one concrete action a data scientist would take.}
\begin{enumerate}
  \item Ask a question: \rule{4in}{0.4pt}
  \item Obtain data: \rule{4in}{0.4pt}
  \item Understand the data: \rule{4in}{0.4pt}
  \item Understand the world: \rule{4in}{0.4pt}
  \item Reports, decisions, solutions: \rule{4in}{0.4pt}
\end{enumerate}
\hint{The lifecycle is a cycle, not a line. Where might this project loop back to an earlier stage?}

\end{document}